\section{Why It Runs Once Per Representation}
\label{sec:ch10-why-it-runs-once-per-representation}

\index{entities@\code{\_entities}!the specification says nothing about how|(}%
Apollo's subgraph specification defines the field this whole chapter is about
in one line:

\begin{graphqlsdl}
_entities(representations: [_Any!]!): [_Entity]!
\end{graphqlsdl}

\index{entities@\code{\_entities}!two normative sentences about the answer}%
and says two things about what comes back. The list
\enquote{must return a list of entity objects that correspond to the provided
representations, in the exact same order}, and
\enquote{entries in the list can be null if no entity exists for a provided
representation}~\autocite{apollosubgraphspec}. Both are about the answer.

\index{entities@\code{\_entities}!nothing normative about resolution}%
There is no third sentence. I read the page looking for one, and there is
nothing anywhere in it about whether a subgraph should answer the list one entry
at a time, all at once, or in any particular order of work. The contract is a
list in, a list out, positionally aligned. How you get from one to the other is
entirely yours.

\index{entities@\code{\_entities}!Hot Chocolate answers it one at a time}%
Hot Chocolate's answer is the one chapter~\ref{ch:data-without-n-plus-1}
already spent a section on, and it is not special to the entity route. A
resolver is a function from one parent to one value. Given three
representations there is one thing the engine can do without knowing anything
about your method, and it does it three times.

\subsection{The Same Mechanism, One Layer Out}

\index{N+1@N+1!chapter 3's shape with a router in it}%
Read that against what chapter~\ref{ch:data-without-n-plus-1} measured and the
two pictures are the same picture. There, the engine resolved \code{sessions},
got four objects, and called the speaker resolver once per object. Here, the
router resolves \code{sessions} against one service, gets four objects, works
out which keys it needs, and calls the entity route of another service with
three of them, whose engine then calls the reference resolver once per entry.
Figure~\ref{fig:ch10-one-fetch-two-costs} is the two costs beside each other.

\begin{figure}[htbp]
  \centering
  % One fetch, two costs. The client asks the router for four sessions and
% each session's speaker name. The router makes two fetches: one statement
% to the Sessions service for the four sessions, then one HTTP request to
% the Speakers service carrying representations for the speakers those
% sessions name. Four sessions name three distinct speakers (Ada Fischer
% speaks twice), and the router itself collapses the four down to three
% before it ever sends the request: the POST to localhost:5002 carries
% three representations, not four. Everything up to and including that one
% POST is identical in both lanes below, and the shared prefix sits at the
% same four x positions in both, the same discipline ch09-two-fetches.tex
% uses for its own shared stages. A brace under lane A names that whole
% prefix explicitly: one request, three representations, deduplicated by
% the router, so a reader cannot mistake what follows for a difference the
% router caused.
%
% What differs is only what the Speakers service does once that request
% lands, and this is where the figure repeats the shape of
% ch03-one-batch.tex with the router standing in for the execution engine
% and a page from another service standing in for the local list ch03
% resolved in memory. Lane A has no DataLoader behind the speaker resolver:
% it is called once per representation, three times, and each call reaches
% SQLite on its own, drawn as three separate ticks rather than an ellipsis,
% exactly as ch03's lane A draws its four. A second brace under those three
% ticks names what N actually is here: one statement per distinct speaker
% on the page, and the page is the Sessions service's. That is the whole
% point of the label, so it has to be the right way round: the Speakers
% service can neither see nor set the length of the list it was handed.
% Lane B sends the resolver through a DataLoader: the same three calls each
% record a key and issue no statement, a dashed batch boundary marks where
% the engine stops asking, and one statement carries all three keys. That
% dashed boundary style is copied unchanged from ch03-one-batch.tex, where
% it carries exactly this meaning: the engine has stopped asking and is
% about to batch what it collected. This figure does not also use it for
% the seam between the router and the Speakers process that
% ch09-two-fetches.tex marks with the same dashed style; that crossing is
% already this book's established idiom from ch09 and is not the argument
% here, and reusing one dashed style for two different meanings in one
% picture would blur the one meaning this figure needs. One more difference
% from ch03 worth naming: there, the DataLoader's batch is what collapses
% four keys to three. Here that collapse already happened one hop upstream,
% inside the router, before either lane's three calls even begin; by the
% time the Speakers service sees the request, in both lanes, there are only
% three keys to resolve, so lane B's batched statement carries three keys
% rather than three deduplicated out of four.
%
% Two stages that an earlier draft of this figure drew are deliberately not
% here, and both were cut for width rather than for taste: the router's
% planning step, which is chapter 8's and chapter 9's subject and not this
% figure's, and a closing response-assembled tick, which ch09-two-fetches.tex
% also does without. The measure is 155mm at this book's geometry (A4, 30mm
% inner, 25mm outer) and every coordinate below is millimetres, so the
% rightmost label has to land inside 155 or the page gains a silent
% overhang. Re-check that if a label here is ever widened.
%
% Same idiom as figures/tikz/ch01-one-field.tex and ch03-one-batch.tex:
% each lane is a timeline read left to right, ticks mark stages, and a
% brace under the stages that matter carries a one-line label. Axis, tick,
% event, span, spanlabel, lane and boundary styles are copied from
% ch09-two-fetches.tex unchanged.
%
% No timings. SPEC decision 19 puts no milliseconds in this book: the
% horizontal axis is sequence only, and no gap below is sized to suggest
% one stage took longer than another. Lane A shows more ticks than lane B
% because it has more distinct stages, not because it is slower; both
% lanes end in the same response to the client.
%
% No quotation marks anywhere in this file. The book's strict Ascii mode
% (see check-chapter.psd1) makes the prose gate read a literal " as a quote
% character, while \" is LaTeX's diaeresis accent and fails to compile
% inside a node. Every identifier below is written bare.
%
% Bare tikzpicture (house style): the figure environment, caption and
% label live at the call site in the section file.
\begin{tikzpicture}[
  x=1mm, y=1mm,
  every node/.append style={font=\sffamily\scriptsize},
  axis/.style={draw, line width=0.4pt,
    -{Straight Barb[length=1.4mm, width=1.6mm]}},
  tick/.style={draw, line width=0.4pt},
  event/.style={anchor=south, align=center, text width=20mm,
    inner sep=1pt},
  span/.style={draw, line width=0.4pt},
  spanlabel/.style={anchor=north, align=center, inner sep=2pt},
  lane/.style={anchor=west, font=\sffamily\scriptsize\itshape},
  boundary/.style={draw, line width=0.4pt, dashed},
]

% ------------------------------------------------------------------- lane A
% At y=20, not 16. The four-line node at x=80 reaches y=15 and this title
% runs to about x=85, so a lower title lands on top of it. ch09-two-fetches.tex
% raises its own lane A title for exactly this reason.
\node[lane] at (0,20)
  {A. one resolver call per representation, each reaching SQLite};

\draw[axis] (0,0) -- (145,0);
\foreach \x in {12,34,56,80,100,116,132}{
  \draw[tick] (\x,-1.5) -- (\x,1.5);
}

\node[event, text width=20mm] at (12,3)
  {client posts to\\localhost:3002\\/graphql};
\node[event, text width=22mm] at (34,3)
  {router posts to\\localhost:5001\\/graphql};
\node[event, text width=20mm] at (56,3)
  {1 statement\\to SQLite:\\4 sessions};
\node[event, text width=26mm] at (80,3)
  {router posts to\\localhost:5002\\/graphql:\\3 representations};
\node[event, text width=14mm] at (100,3)  {rep 1\\stmt 2};
\node[event, text width=14mm] at (116,3)  {rep 2\\stmt 3};
\node[event, text width=14mm] at (132,3)  {rep 3\\stmt 4};

\draw[span] (12,-4) -- (12,-7) -- (80,-7) -- (80,-4);
\node[spanlabel, text width=58mm] at (46,-8)
  {identical in both lanes: one POST to\\localhost:5002, carrying 3
   representations,\\deduplicated from 4 by the router};

\draw[span] (94,-4) -- (94,-7) -- (138,-7) -- (138,-4);
\node[spanlabel, text width=42mm] at (116,-8)
  {3 statements: one per distinct\\speaker on the page, and the\\page is
   the Sessions service's};

% ------------------------------------------------------------------- lane B
\begin{scope}[shift={(0,-48)}]
  \node[lane] at (0,20)
    {B. the same three calls, one batched statement};

  \draw[axis] (0,0) -- (145,0);
  \foreach \x in {12,34,56,80,134}{
    \draw[tick] (\x,-1.5) -- (\x,1.5);
  }
  \foreach \x in {98,103,108}{
    \draw[tick] (\x,-1.5) -- (\x,1.5);
  }
  \draw[boundary] (120,-3) -- (120,6);

  \node[event, text width=20mm] at (12,3)
    {client posts to\\localhost:3002\\/graphql};
  \node[event, text width=22mm] at (34,3)
    {router posts to\\localhost:5001\\/graphql};
  \node[event, text width=20mm] at (56,3)
    {1 statement\\to SQLite:\\4 sessions};
  \node[event, text width=26mm] at (80,3)
    {router posts to\\localhost:5002\\/graphql:\\3 representations};
  % Raised to y=7 so that its box clears the four-line node at x=80, and
  % centred at 107 rather than over its own tick cluster so that 24mm of
  % width fits without the two boxes touching. 22mm would fit over the
  % cluster and hyphenates the word resolver, which reads badly.
  \node[event, text width=24mm] at (107,7)
    {speaker resolver x3:\\records key, no query};
  \node[event, text width=18mm] at (134,3)
    {1 statement:\\3 keys};

  % One label below this lane's axis and no more. An earlier layout put
  % three there, and at this spacing their boxes overlapped into an
  % unreadable run: the three stages lane B annotates sit inside 40mm,
  % where ch03-one-batch.tex had 70mm for two. The two labels dropped said
  % what the event nodes above the axis already say.
  \node[spanlabel, text width=28mm] at (120,-8)
    {batch boundary:\\engine stops asking};
\end{scope}

\end{tikzpicture}

  \caption{The same entity fetch, answered two ways. Everything up to the
  fourth tick is identical: the router gets four sessions from the Sessions
  service in one statement, then posts one request to the Speakers service
  carrying three representations, having dropped the duplicate itself. What
  that service does with the list after it lands is the whole of the
  difference, and lane A's three statements are one per distinct speaker on a
  page the Sessions service chose. This is
  figure~\ref{fig:ch03-one-batch} with a router where the execution engine was.}
  \label{fig:ch10-one-fetch-two-costs}
\end{figure}

\index{N+1@N+1!the unit of repetition changes}%
What changed between the two chapters is the unit and who controls it. In
chapter~\ref{ch:data-without-n-plus-1} the repeated statement went to a database
in the same process, over a list that same process had just produced, and
\code{[UsePaging]} on that field capped it. Here the list arrives from a service
this one does not own, and the cap on its length is a paging decision made in
the Sessions codebase. A subgraph author cannot bound N by reading their own
code.

\index{N+1@N+1!what the wire hides}%
There is a second difference and it runs the other way, which is the only
reason this is not worse than chapter~\ref{ch:data-without-n-plus-1}'s version.
The repetition stops at the process boundary. One HTTP request carried the
whole list, so the network cost is one round trip whatever the resolver does
with it, and the fan-out lands entirely on the database behind the subgraph.
That is genuinely better than a gateway that made one call per row, and it is
also exactly why nobody notices: the expensive thing is now one layer below the
place anyone is watching.

\subsection{Two Things That Do Not Make It Better}

\index{N+1@N+1!not fixed by the router}%
The router will not save you, and it is worth being precise about why rather
than hoping. It de-duplicates, as section~\ref{sec:ch10-three-representations-not-four}
showed, and that is the whole of what it can do from outside. Collapsing three
lookups into one query requires knowing that all three go to the same table,
that the table can be queried by a set of keys, and that the results can be put
back in order. None of that is visible over HTTP. The router has a list of
opaque identifiers and a schema.

\index{N+1@N+1!not fixed by the entity route being one field}%
And the fact that \code{\_entities} is a single field with a list argument, which
looks like a batching API and is one, does not batch anything by itself. It gets
the list to your door. What happens after that is a property of the code you
wrote, which is the whole reason this chapter exists rather than being a
paragraph in chapter~\ref{ch:first-subgraph}.
\index{entities@\code{\_entities}!the specification says nothing about how|)}%
